%======================================================================
% 6. VERIFICATION AGAINST THE TRUE SIMULATOR
%======================================================================
\section{Verification against the true simulator}
\label{sec:verification}

Every result so far rests on the surrogate moment engine, so five designs
were re-examined against the trajectory simulator itself: the deterministic
and robust knees, the robust matched design of Table~\ref{tab:knees}, the
most fragile deterministic design (the front's worst shoulder-heating
exceedance probability --- the audit headline), and the minimum-dispersion
robust design (the floor claim). Each was propagated with $N=300$
uncertainty samples drawn from the frozen distributions
($t_{\max}=\SI{14400}{\s}$), 1{,}500 simulator runs in total, about 23~s
each. There were no integration failures and 1{,}499 of 1{,}500 runs
captured.

Table~\ref{tab:verify} summarizes the comparison for the three designs that
stress the claims. Expected heat load is verified to within $-3.2$ to
$+1.4\%$ across all five designs, and $\sigma_{Q_s}$ to within a few
percent (MC slightly lower throughout); in particular the floor survives:
the minimum-dispersion design measures $\sigma_{Q_s}=
\SI{15.4}{\mega\joule\per\square\metre}$ in simulation against the claimed
\SI{16.4}{\mega\joule\per\square\metre} --- the floor is real, and the
inflated-$\sigma$ construction errs on the conservative side. The
peak-heating $\mathrm{E}+2\sigma$ constraints are met with 1--3\%
margin. The audit's headline number is confirmed at its worst point: the
fragile deterministic design shows $P(q_{\mathrm{shldr}}>\text{limit})=
29/299=9.7\%$ in simulation against the surrogate's promised 10.0\%, and
$P(q_{\mathrm{stag}}>\text{limit})=7.4\%$ against 6.8\%. The 299
denominators are exact bookkeeping rather than a quiet exclusion: the
fragile design's shallowest draw ($\gamma_0=-0.52^\circ$, three standard
deviations above nominal, with $k_{CL}=1.07$) lofts into a long skip
trajectory that is still airborne at the integration horizon
$t_{\max}=\SI{14400}{\s}$ --- a physical escape of the terminal event
rather than an integration failure, and non-capturing under both
guidance laws. Its recorded heat load and shoulder heating at cutoff lie
far below the limits; counting that run as an exceedance or as a
non-exceedance brackets the estimate between 29/300 and 30/300
(9.7--10.0\%), so the verification statement is insensitive to its
treatment. (Under the alternate law of Section~\ref{sec:guidance} the
same brackets are 17.0--17.3\% around the reported 17.1\%.) Both knees
show zero exceedances in 300 runs, which bounds their true exceedance
probabilities below $\approx$1\% (rule of three). Slow-family
probabilities agree within about six percentage points.

\begin{table}[htbp]
\centering
\caption{Monte-Carlo verification of the surrogate moment engine
  ($N=300$ true-simulator runs per design). Diff. = (MC $-$ surrogate)
  in \% of the surrogate value.}
\label{tab:verify}
\begin{tabular}{lccc}
\toprule
quantity & det knee & robust knee & det fragile \\
\midrule
$\mathrm{E}[Q_s]$ diff.        & $+0.1\%$  & $+1.4\%$  & $-3.2\%$ \\
$\sigma_{Q_s}$ MC / predicted  & 23.9/24.8 & 36.2/37.3 & 65.9/66.7 \\
$\mathrm{E}[sg]$ diff.         & $-24.6\%$ & $-8.8\%$  & $-20.1\%$ \\
$P(\mathrm{slow})$ MC / pred.  & 0.40/0.45 & 0.46/0.48 & 0.49/0.49 \\
$P(q_{\mathrm{stag}})$ MC / pred. & 0/300 / $<0.1\%$ & 0/300 & 7.4\%/6.8\% \\
$P(q_{\mathrm{shldr}})$ MC / pred. & 0/300 / $<0.1\%$ & 0/300 & \textbf{9.7\%/10.0\%} \\
\bottomrule
\end{tabular}
\end{table}

One systematic discrepancy survives verification and is disclosed rather
than tuned away: predicted $\mathrm{E}[sg]$ is high by 8.8--28.1\% on all
five designs (standard errors 2.7--10.1\%). The bias has a single,
identifiable cause. The mixture prediction averages a lognormal for each
family, and the mean of a lognormal amplifies the fitted variance; for the
slow family, whose fitted log-standard-deviation is large by necessity
(Section~\ref{sec:gp}), that amplification is severe --- the same mechanism
that produced the 3.5$\times$ nominal-range extrapolation of
Section~\ref{sec:exploit}, here in its mild, moment-level form. Monte-Carlo
truth for the five designs puts $\mathrm{E}[sg]$ at 4.1--5.1~Mm. What
this bias does to the matched-pair comparison is a fair question, and
the answer is not that it cancels: the window holds $\mathrm{E}[sg]$
equal only in surrogate units, so what matters is the
\emph{differential} bias between matched designs. The measured bias is
design-specific --- multiplicative factors 1.10--1.39 across the five
verified designs, with no systematic difference between the fronts
(deterministic 1.25--1.33, robust 1.10--1.39) --- and at the
matched knee pair verified end to end (deterministic knee against the
robust design of Table~\ref{tab:knees}) it moves the range
ratio from $-2.1\%$ in surrogate units to $-6.6\%$ in simulation, while
the comparison itself survives: the heat-load leg of the window holds
with margin ($-6.4\%$), and the $\sigma_{Q_s}$ cut grows from 3.7\% to
5.1\%.

The differential bias can also be measured directly, at the pairs where
it had the best chance to distort the comparison. The five most
adversarial matched pairs were re-verified end to end against the
simulator: one pair per expected-heat-load quintile, each deterministic
member chosen with a predicted $\sigma_{Q_s}$ cut at least the
pair-level median of 27.5\% and the least range slack in its window ---
precisely where a surrogate error could most plausibly manufacture or
break a match. Both members of every pair were propagated with the same
$N=300$ shared uncertainty samples (3{,}000 further simulator runs;
2{,}999 captured, the single non-capture being one shallow lofted draw
of the fifth pair's deterministic member, of the kind already described
above). The comparison survives intact. The predicted $\sigma_{Q_s}$
cuts of $+30.9$ to $+67.1\%$ measure $+32.0$ to $+68.5\%$ in simulation
--- per-pair agreement within 1.5 percentage points, and the true cut
the larger in four of five pairs. The heat-load leg of the window is
likewise confirmed: member ratios of expected heat load $0.42$--$0.74$
predicted against $0.40$--$0.74$ measured, the largest deviation
(pair~4: 0.402 against 0.420) understating the robust advantage. The
range ratios move from the matched $-2.0\%$ in surrogate units to
$-5.2$ to $+8.1\%$ measured (standard errors near $\pm 5\%$): the
differential bias is real but bounded, and at these five pairs it runs
in the robust front's favour --- the robust members' measured
$\mathrm{E}[sg]$ inflation is smaller on average (1.23 against 1.26;
factors 1.14--1.37 deterministic against 1.10--1.33 robust, all inside
the envelope above) --- so the mean measured range ratio is 1.01:
matched robust designs give up no range in truth. The matched-pair statistics
of Section~\ref{sec:matched} therefore rest on direct measurement, not
only on the bounds below.

Two bounds cover the pairs not verified directly. First, widening or
removing the range leg of the window leaves the pair statistics intact
(Section~\ref{sec:sens}: preference 97.2--98.1\%, median cut
27.5--32.2\%), so no differential bias of plausible magnitude can
overturn the comparison through the matching. Second, the Monte-Carlo
standard error on $\mathrm{E}[sg]$ at $N=300$ is itself 3.5--4\%, so a
2\% range window is beyond the resolution of any verification budget
this study could afford; the heat-load leg, verified to $-3.2$ to
$+1.4\%$, is what carries the matching. Absolute range values quoted
from the surrogate should be read as conservative-high by 9--37\%
across the fifteen verified designs.

\subsection{Guidance-law sensitivity}
\label{sec:guidance}

The guidance realization was the dominant residual in the end-to-end
validation (Appendix~\ref{app:variants}), so the verification was repeated
under the alternate bank law of Table~\ref{tab:variants} --- Eq.~2.40 as
printed with the 82$^\circ$ command cap, the variant that reproduces the
reference track length to $-1.1\%$ --- re-running the same five designs at
the same 300 uncertainty samples (1{,}500 further simulator runs). The
shifts are systematic rather than erratic: expected heat load moves $-1$
to $-7\%$, expected range $-5$ to $-13\%$, and heat-load dispersion
\emph{increases} by 7--24\% under the more aggressive modulation. Two
results matter. The verified knees remain safe under both laws (zero
exceedances in 300 runs at the deterministic knee and the matched design;
2/300 shoulder exceedances at the robust knee). And the hidden fragility
of the deterministic front survives and amplifies: at the fragile design,
shoulder-heating exceedances rise from 29/299 (9.7\%) to 51/299 (17.1\%)
and stagnation-heating exceedances from 22/299 to 42/299. The concealed
risk is thus not an artefact of one guidance realization --- if anything,
the alternate law is harsher on the point-feasible design. Absolute
dispersion levels, including the floor's magnitude (15.4 to
\SI{19.1}{\mega\joule\per\square\metre} at the floor design), are
guidance-dependent exactly as the claims accounting of
Section~\ref{sec:validation} anticipated; the comparison conclusions, and
the floor's geometry-independence, are not.
